\chapter{Маршрутизация}
\label{ch:chap4}

\begin{figure}[ht!]
    \centering
    \includegraphics[width=1\textwidth]{screenshots/4_1_1.jpg}
    \caption{Топология лабораторной работы по маршрутизации}
    \label{fig:4_1_1}
\end{figure}
\clearpage
\begin{figure}[ht!]
    \centering
    \includegraphics[width=1\textwidth]{screenshots/4_1_2.jpg}
    \caption{Назначение ip-адресов на маршрутизаторе S\_R2}
    \label{fig:4_1_2}
\end{figure}
\clearpage
\begin{figure}[ht!]
    \centering
    \includegraphics[width=1\textwidth]{screenshots/4_1_3.jpg}
    \caption{Назначение ip-адресов на маршрутизаторе S\_R1}
    \label{fig:4_1_3}
\end{figure}
\clearpage
\begin{figure}[ht!]
    \centering
    \includegraphics[width=1\textwidth]{screenshots/4_1_4.jpg}
    \caption{Назначение ip-адресов на маршрутизаторе S\_R3}
    \label{fig:4_1_4}
\end{figure}
\clearpage
\begin{figure}[ht!]
    \centering
    \includegraphics[width=1\textwidth]{screenshots/4_1_5.jpg}
    \caption{Назначение ip-адресов на маршрутизаторе устройства провайдера}
    \label{fig:4_1_5}
\end{figure}
\clearpage
\begin{figure}[ht!]
    \centering
    \includegraphics[width=1\textwidth]{screenshots/4_1_6.jpg}
    \caption{Назначение ip-адресов на VPC4}
    \label{fig:4_1_6}
\end{figure}
\clearpage
\begin{figure}[ht!]
    \centering
    \includegraphics[width=1\textwidth]{screenshots/4_1_7.jpg}
    \caption{Назначение ip-адресов на VPC5}
    \label{fig:4_1_7}
\end{figure}
\clearpage
\begin{figure}[ht!]
    \centering
    \includegraphics[width=1\textwidth]{screenshots/4_1_8.jpg}
    \caption{Настройка роутинг на маршрутизаторе S\_R1}
    \label{fig:4_1_8}
\end{figure}
\clearpage
\begin{figure}[ht!]
    \centering
    \includegraphics[width=1\textwidth]{screenshots/4_1_9.jpg}
    \caption{Настроенный роутинг на маршрутизаторе S\_R2}
    \label{fig:4_1_9}
\end{figure}
\clearpage
\begin{figure}[ht!]
    \centering
    \includegraphics[width=1\textwidth]{screenshots/4_1_10.jpg}
    \caption{Настройка роутинг на маршрутизаторе S\_R3}
    \label{fig:4_1_10}
\end{figure}
\clearpage
\begin{figure}[ht!]
    \centering
    \includegraphics[width=1\textwidth]{screenshots/4_1_11.jpg}
    \caption{Проверка доступности VPC5 из VPC4 через ping}
    \label{fig:4_1_11}
\end{figure}
\clearpage
\begin{figure}[ht!]
    \centering
    \includegraphics[width=1\textwidth]{screenshots/4_1_12.jpg}
    \caption{Проверка доступности VPC4 из VPC5 через ping}
    \label{fig:4_1_12}
\end{figure}
\clearpage
\begin{figure}[ht!]
    \centering
    \includegraphics[width=1\textwidth]{screenshots/4_1_13.jpg}
    \caption{Проверка доступности подсети 192.168.2.0/30 из VPC4 через ping}
    \label{fig:4_1_13}
\end{figure}
Не получилось, потому что в настройках маршрутизатора S\_R1 не был указан маршрут по умолчанию на подсеть 192.168.2.0/30. После добавления маршрута по умолчанию на S\_R2 был сконфигурирован доступ только к подсети 192.168.1.0/24, где находится VPC5. Доступ к искомой подсети не был сконфигурирован.
\clearpage
\begin{figure}[ht!]
    \centering
    \includegraphics[width=1\textwidth]{screenshots/4_1_14.jpg}
    \caption{Настройка доступа к нужной подсети 192.168.2.0/30 на маршрутизаторе s\_r1}
    \label{fig:4_1_14}
\end{figure}
\clearpage
\begin{figure}[ht!]
    \centering
    \includegraphics[width=1\textwidth]{screenshots/4_1_15.jpg}
    \caption{Успешная доступность нужной подсети 192.168.2.0/30 из VPC4 после настройки маршрута на S\_R1}
    \label{fig:4_1_15}
\end{figure}
